\documentclass[10pt,a4paper]{amsart}
\usepackage[margin=19mm]{geometry}
\usepackage{amsmath,amssymb,mathtools}
\usepackage[scr=rsfs,cal=euler]{mathalfa}
\usepackage{xfrac}
\usepackage[unicode,hidelinks,bookmarks=false]{hyperref}
\newcommand{\ie}{i.e.\ }
\newcommand{\cf}{cf.\ }
\newcommand{\resp}{respectively\ }
\newcommand{\Subsection}{\subsection}
\providecommand{\nobreakdash}{\nobreak\hbox{-}}
\setlength{\emergencystretch}{2em}
\multlinegap0pt
\usepackage{amssymb}
\usepackage{bm}
\usepackage[shortlabels]{enumitem}
\usepackage{mathtools}
\usepackage{xcolor}
\usepackage{tikz}
\usepackage{float}
\usepackage{algorithm}\usepackage{algpseudocode}
\newtheorem{proposition}{Proposition}
\newtheorem{lemma}{Lemma}
\usepackage{cleveref}
\crefname{proposition}{Proposition}{Propositions}
\crefname{lemma}{Lemma}{Lemmas}
\newcommand{\D}{\Delta}
\newcommand{\Dx}{{\Delta x}}
\newcommand{\beq}{\begin{equation}}
\newcommand{\eeq}{\end{equation}}
\newcommand{\g}{\gamma}
\newcommand{\lc}{\left(}
\newcommand{\p}{\psi}
\newcommand{\rc}{\right)}
\newcommand{\ux}{{\underline x}}
\title{Finite difference methods for a continuous-time heterogeneous agent model with recursive utility}
\author{Yves Achdou, Qing Tang}
\date{Source: arXiv:2606.23408v1 (CC BY 4.0)}
\begin{document}
\maketitle

\noindent\emph{Source and licence.} Finite difference methods for a continuous-time heterogeneous agent model with recursive utility. Version arXiv:2606.23408v1 (CC BY 4.0), distributed by arXiv under the Creative Commons Attribution 4.0 International licence, http://creativecommons.org/licenses/by/4.0/, as recorded in the arXiv licence metadata for this exact version and shown on the abstract page https://arxiv.org/abs/2606.23408v1. Full licence notice: \texttt{licenses/arxiv-2606.23408-CC-BY-4.0.txt}.\par\smallskip

\noindent\textit{Editorial scope.} Counted unit: the article's lemma FD (source0.tex lines 1132--1144). The article's complete original proof of it is delivered with it (source0.tex lines 1145--1194, one \texttt{proof} environment of 50 lines). No mathematics is written, repaired or re-derived: the statement and the proof are the article's own lines, in the article's own notation, and no retained line is substituted or re-flowed.  Retained uncounted as the source-local context the proof needs: lines 344--346; lines 348--356; lines 371--383; lines 386--396.  Nothing was omitted from any retained span, so the package carries no omission record.  No retained line contains a citation, so the wrapper carries no bibliography and no bibliography entry.  No retained line contains a figure environment, an image-inclusion command or drawing code, so no image is delivered and the package declares no asset dependency.  Editorial formatting only, in addition: an AMS \texttt{amsart} preamble that restates the article's own declarations and macros used by the retained lines, namely the environments \texttt{proposition}, \texttt{lemma} on separate counters, together with the standard packages those lines need.  The retained environments print in this document as Proposition 1, Lemma 1 in order of appearance, whereas the article prints the same items with the numbering of its own full document.  The complete native source of the article as distributed in the arXiv e-print for this exact version is bundled unchanged as \texttt{sources/203429/source0.tex}.
\begin{proposition}\label{Prop: v sol}
Assume $\theta \geq 1$. There exists a unique viscosity solution $(v_1,v_2)$ which is $C^1$ and strictly concave. Moreover, $v_j\in W^{2,\infty}_{loc}(\ux,+\infty)$ and $v_2\geq v_1$. 
\end{proposition}

\begin{proposition}\label{Prop: saving}
Assume $\theta \geq 1$. The optimal saving policy $s_1$ has the following properties: $s_1(x)< 0$ for all $x> \ux$ and $s_1(\ux)=0$. If furthermore
\beq\label{rho-r2}
\begin{aligned}
\lc \rho-r\rc(r\ux+y_2)^{-1/\psi}+\lambda_2\lc (r\ux+y_2)^{-1/\psi}-(r\ux+y_1)^{-1/\psi}\rc< 0,
\end{aligned}
\eeq
then $s_2(\ux)>0$. Moreover, $Dv_1(\ux)>Dv_2(\ux)$. 
\end{proposition}

\begin{proposition}\label{Prop: Dv1 ux}
We make the same assumptions as in \cref{Prop: saving}. Near $x=\ux$, the functions $s_1$, $D^2v_1$ and $Dv_1$ have the following behavior: 
\beq\label{Eq: s1D2v1 ux}
s_1(x)D^2v_1(x)= \varkappa+o(1),
\eeq
\beq\label{Eq: Dv1 ux}
Dv_1(x)\sim \rho\frac{((1-\g)v_1(\ux))^{\frac{\p^{-1}-\g}{1-\g}}}{ (r\ux+y_1)^{\p^{-1}}}-\sqrt{\frac{2\varkappa(x-\ux)}{\p \lc r\ux+y_1\rc^{1+\p^{-1}}} ((1-\g)v_1(\ux))^{\frac{\p^{-1}-\g}{1-\g}}},
\eeq
where $\varkappa$ is given by
\beq\label{Eq: varkappa}
\varkappa:=\lc \frac{\rho}{\theta}-r\rc Dv_1(\ux)+\lambda_1(Dv_1(\ux)-Dv_2(\ux))-H_v(\ux,y_1,v_1(\ux),Dv_1(\ux))Dv_1(\ux).
\eeq
\end{proposition}

\begin{lemma}\label{Prop: D2v1 ux}
We make the same assumptions as in \cref{Prop: saving}. There exists $\eta\in (\ux,\bar{x})$ such that for all $x\in (\ux,\eta)$, 
\beq\label{Eq: -D2v1 bound}
0<-Dv^2_1(x)\leq \sqrt{\frac{\varkappa}{\p \lc r\ux+y_1\rc^{1+\p^{-1}}(x-\ux)}\cdot ((1-\g)v_1(\ux))^{\frac{\p^{-1}-\g}{1-\g}}}.
\eeq
and 
\beq\label{Eq: -s1 bound}
0<-s_1(x)<2\sqrt{\varkappa \p \lc r\ux+y_1\rc^{1+\p^{-1}}((1-\g)v_1(\ux))^{\frac{\p^{-1}-\g}{\g-1}}(x-\ux)}
\eeq

\end{lemma}

\begin{lemma}\label{Lemma: error FD}
We make the same assumptions as in \cref{Prop: saving}, then there exists a constant $C$ (uniform in $\D x$) such that
\beq\label{Eq: FD erro v2}
\max\lc\max_{\,x\in \mathcal{G}^{\D x},\,x>\ux}\left\{\D^-{v_{2}(x)}-Dv_{2}(x)\right\},\,\max_{\,x\in \mathcal{G}^{\D x},\,x<\bar{x}}\left\{Dv_{2}(x)-\D^+{v_{2}(x)}\right\}\rc\leq C\D x.
\eeq
There exists $\eta$ sufficiently small (uniformly in $\D x$) such that
\beq\label{Eq: FD boundDx}
\max_{\,x\in \mathcal{G}^{\D x},\,\ux<x<\eta}\left\{\D^-{v_{1}(x)}-Dv_{1}(x)\right\}\leq C\sqrt{\D x},
\eeq
\beq\label{Eq: s1FD boundDx}
\max_{\,x\in \mathcal{G}^{\D x},\,\ux<x<\eta}\left\{-s_1(x)\lc\D^-{v_{1}(x)}-Dv_{1}(x)\rc\right\}\leq C\D x.
\eeq
\end{lemma}

\begin{proof}
In the proof, the constant $C$ may change from line to line but is always independent of $\D x$. From \cref{Prop: v sol}, we know that $v_j\in W^{2,\infty}_{loc}(\ux,\bar{x}]$. Since $s_2(\ux)>0$,  $D^2 v_2(\ux)$ is bounded. This implies $D^2 v_2\in L^\infty(\ux,\bar{x})$ and there exists $C_2$ (independent of $\D x$) such that 
$$
\begin{aligned}
\max\lc\max_{x\in \mathcal{G}^{\D x},\,x>\ux}\left\{\D^-{v_{2}(x)}-Dv_{2}(x)\right\},\,\max_{x\in \mathcal{G}^{\D x},\,x<\bar{x}}\left\{Dv_{2}(x)-\D^+{v_{2}(x)}\right\}\rc\leq C_2\Dx.
%\max\lc\sup_{x>\ux}\left\{\D^-{v_{1}(x)}-Dv_{1}(x)\right\},\,\sup_{x<\bar{x}}\left\{Dv_{1}(x)-\D^+{v_{1}(x)}\right\}\rc\leq C_1\Dx.
\end{aligned}
$$
To prove \cref{Eq: FD boundDx} and \cref{Eq: s1FD boundDx}, we consider two cases: $x=\ux+\D x$ and $\ux+\D x<x<\eta$.\par
Case 1:  $x=\ux+\D x$. From the strict concavity of $v_1$, 
$$
0<\D^-{v_{1}(x+\D x)}-Dv_{1}(x+\D x)<Dv_{1}(x)-Dv_{1}(x+\D x)\leq C\sqrt{\D x},
$$ 
where for the last inequality we used \cref{Eq: Dv1 ux}. From \cref{Eq: -s1 bound}, we know $0<-s_1(x+\D x)\leq C\sqrt{\D x}$, hence
$
-s_1(x+\D x)\lc\D^-{v_{1}(x+\D x)}-Dv_{1}(x+\D x)\rc \leq C\D x.
$\par
 Case 2: $\ux+2\D x\leq x< \eta$. By using a Taylor expansion with integral remainder, we deduce
$$
\D^-{v_{1}(x)}-Dv_{1}(x)=-\frac{1}{\D x}\int_{x-\D x}^{x}(z-(x-\D x))D^2v_1(z)dz.
$$
From $x-\D x\leq z<\eta$ and \cref{Prop: D2v1 ux}, we deduce that $0<-D^2v_1(z)<\frac{C}{\sqrt{z-\ux}}\leq \frac{C}{\sqrt{x-\D x-\ux}}\leq \frac{C}{\sqrt{\D x}}$, where for the last inequality we used $x\geq \ux+2\D x$. Hence,  
$$
\D^-{v_{1}(x)}-Dv_{1}(x)\leq \frac{C}{(\D x)^{1/2}}\frac{1}{\D x}\int_{x-\D x}^{x}(z-(x-\D x))dz\leq C(\D x)^{1/2}.
$$
Since $0<-s_1(x)\leq C\sqrt{x-\ux}$ for $\ux+2\D x\leq x< \eta$, 
$$
\begin{aligned}
-s_1(x)\lc\D^-{v_{1}(x)}-Dv_{1}(x)\rc={}&\frac{-s_1(x)}{\D x}\int_{x-\D x}^{x}(z-(x-\D x))\lc-D^2v_1(z)\rc dz\\
\leq {}&\frac{C\sqrt{x-\ux}}{\sqrt{x-\D x-\ux}}\frac{1}{\D x}\int_{x-\D x}^{x}(z-(x-\D x))dz\\
\leq {}&\frac{C(\sqrt{x-\D x-\ux}+\sqrt{\D x})}{\sqrt{x-\D x-\ux}}\D x\leq C\D x.
\end{aligned}
$$
For the last inequality, we observe that $\frac{\sqrt{\D x}}{\sqrt{x-\D x-\ux}}\leq 1$ since $x\geq \ux+2\D x$. 
%This yields, for sufficiently small $\D x$, 
%$$
%\begin{aligned}
%Dv_{1}(\ux)-\D^+{v_{1}(\ux)}={}& \frac{1}{\D x}\int_{\ux}^{\ux+\D x}\left[Dv_1(\ux)-Dv_{1}(x) \right]dx\\
%\leq {}& \sqrt{\frac{3\varkappa \D x}{\p \lc r\ux+y_1\rc^{1+\p^{-1}}}((1-\g)v_1(\ux))^{\frac{\p^{-1}-\g}{1-\g}}}.
%\end{aligned}
%$$

%$$
%D^2v_1(x)\sim \sqrt{\frac{\varkappa}{2\p \lc r\ux+y_1\rc^{1+\p^{-1}}(x-\ux)}\cdot ((1-\g)v_1(\ux))^{\frac{\p^{-1}-\g}{1-\g}}}.
%$$
%This yields
%$$
%Dv_{1}(\ux)-\D^+{v_{1}(\ux)} \leq \sqrt{\frac{\varkappa \D x}{\p \lc r\ux+y_1\rc^{1+\p^{-1}}}((1-\g)v_1(\ux))^{\frac{\p^{-1}-\g}{1-\g}}}
%$$
\end{proof}

\end{document}
